\label{sec:introduccion}
Aproximar fenómenos físicos gobernados por ecuaciones diferenciales exige decidir cómo distribuir recursos computacionales. Una solución puede tener buen error promedio y aún fallar cerca de una interfaz o un frente. La precisión local importa tanto como la capacidad de reconstruir un campo completo; su evaluación requiere referencias y presupuestos comparables \cite{grossmann2024,krishnapriyan2021}.

Las redes neuronales informadas por física combinan aproximación y restricciones diferenciales; las PIKANs utilizan funciones univariadas aprendibles como componentes del aproximador \cite{raissi2019,liu2024kan}. Su flexibilidad abre posibilidades para adaptar representación durante el entrenamiento, junto a estrategias existentes de muestreo, ponderación y descomposición. Esta tesis estudia la decisión de reasignación, apoyándose en esos antecedentes.

La propuesta investiga si un controlador puede asignar resolución con presupuesto fijo, conservar lo aprendido durante cambios geométricos y evitar trasladar error entre regiones o campos. La contribución se evaluará mediante ablaciones y costo total, no mediante la presunción de que adaptar una grilla sea nuevo. Los fundamentos se desarrollan en el capítulo 6; el capítulo 7 fija la evaluación antes de implementar. El caso piezoeléctrico constituye una prueba de extensión a campos acoplados.
